% Qubits desde cero · 03 -- ¿Cómo describimos un qubit?
% Tema: Signature (serie principal)
% Build:  ./build.sh qubits-desde-cero-03.tex   (o en Overleaf con pdfLaTeX)
\documentclass[12pt]{beamer}
\usetheme{igsignature}

\ighandle{@quantum_danesh}
\igseries{Qubits desde cero}
\igissue{03}

% --- Diagramas de Stern-Gerlach (los mismos del post 02) -------------------
% Un imán: \pic at (x,y) {sg={eje/ángulo de la flecha/color de la flecha}};
\tikzset{pics/sg/.style args={#1/#2/#3}{code={
  \draw[igAccentB, line width=0.7pt, fill=igPanel, rounded corners=1.6mm] (-7,-4.6) rectangle (7,4.6);
  \draw[igMuted!60, line width=0.4pt] (-3.4,0) circle (2.6);
  \draw[#3, line width=0.9pt, -{Latex[length=1.3mm]}] (-3.4,0) ++(#2+180:1.6) -- ++(#2:3.4);
  \node[text=igFG, font=\footnotesize] at (3.2,2.1) {$+#1$};
  \node[text=igFG, font=\footnotesize] at (3.2,-2.1) {$-#1$};
  \node[text=igMuted, font=\small] at (0,-7.5) {SG$_{#1}$};}}}
\newcommand\sgpass[2]{% de la salida + de un imán a la entrada del siguiente
  \draw[igAccentB, line width=0.8pt] (#1+7,2.1) .. controls +(4,0) and +(-4,0) .. (#2-7,0);}
\newcommand\sgblock[1]{% salida - bloqueada
  \draw[igMuted, line width=0.8pt] (#1+7,-2.1) -- (#1+10,-2.1);
  \draw[igMuted, line width=1.6pt] (#1+10,-3.6) -- (#1+10,-0.6);}
\newcommand\sgout[3]{% las dos salidas del último imán, con sus etiquetas
  \draw[igAccentB, line width=0.8pt] (#1+7,2.1) -- (#1+10,4.5);
  \draw[igAccentB, line width=0.8pt] (#1+7,-2.1) -- (#1+10,-4.5);
  \node[anchor=west, fill=igHL, draw=igTrack, rounded corners=1.2mm, text=igAccent,
    font=\small, inner sep=1mm] at (#1+10,4.5) {#2};
  \node[anchor=west, fill=igHL, draw=igTrack, rounded corners=1.2mm, text=igAccent,
    font=\small, inner sep=1mm] at (#1+10,-4.5) {#3};}
% etiqueta de fila, en gris, a la derecha de un vector columna
\newcommand\rowlabel[1]{\text{\small\color{igMuted}$\leftarrow$ #1}}

\begin{document}

\coverslide{Qubits desde cero}
  {¿Cómo describimos un qubit?}
  {Del imán de Stern--Gerlach a un vector de dos números.}

% ---------------------------------------------------------------------------
\begin{frame}{Sistema y estado}{Dos ideas distintas}
\begin{center}
\begin{tikzpicture}[x=1mm,y=1mm,scale=1.2,transform shape]
  \node[draw=igMuted, fill=igPanel, rounded corners=1mm, minimum width=11mm,
        minimum height=7mm, text=igMuted, font=\footnotesize] at (0,0) {horno};
  \draw[igAccentB, line width=0.8pt] (5.5,0) -- (25,0);
  \foreach \x in {9,13,17,21} \fill[igFG] (\x,0) circle (0.45);
  \pic at (32,0) {sg={z/90/igAccentB}};
  \sgout{32}{?}{?}
\end{tikzpicture}
\end{center}
\begin{defbox}[Sistema]
Lo que medimos: el espín de un átomo de plata. Siempre es el mismo y el
imán solo puede dar \key{dos resultados}.
\end{defbox}
\begin{thmbox}[Estado]
Cómo lo hemos preparado: qué filtros atravesó. Cambia de un experimento a
otro y decide las \hl{probabilidades}.
\end{thmbox}
\end{frame}

% ---------------------------------------------------------------------------
\begin{frame}{Dos salidas, dos etiquetas}{Del imán a la notación}
\begin{center}
\begin{tikzpicture}[x=1mm,y=1mm,scale=1.45,transform shape]
  \useasboundingbox (-1,-9) rectangle (52,7);
  \draw[igAccentB, line width=0.8pt] (0,0) -- (21,0);
  \pic at (28,0) {sg={z/90/igAccentB}};
  \sgout{28}{\alt<2->{$\ket{0}$}{$+z$}}{\alt<2->{$\ket{1}$}{$-z$}}
\end{tikzpicture}
\end{center}
El analizador en $z$ solo tiene dos salidas: arriba ($+z$) o abajo ($-z$).

\uncover<2->{Les ponemos nombre:
\[
  +z \;\longrightarrow\; \ket{0},
  \qquad
  -z \;\longrightarrow\; \ket{1}.
\]
Elegir el analizador es elegir la \hl{base} en la que describimos el qubit.}
\end{frame}

% ---------------------------------------------------------------------------
\begin{frame}{Cada etiqueta es un vector}{La base}
{\large
\[
  \ket{0} = \begin{pmatrix} 1 \\ 0 \end{pmatrix},
  \qquad
  \ket{1} = \begin{pmatrix} 0 \\ 1 \end{pmatrix}
\]}%
Un \key{ket} $\ket{\,\cdot\,}$ es un vector columna de dos componentes.

$\ket{0}$ lo tiene todo en la primera componente, la de la salida $+z$;
$\ket{1}$, en la segunda, la de la salida $-z$.
\end{frame}

% ---------------------------------------------------------------------------
\begin{frame}{Cualquier estado}{El ket}
{\large
\begin{align*}
  \ket{\psi} &= \alpha\ket{0} + \beta\ket{1}\\
  &\uncover<2->{= \alpha\begin{pmatrix} 1 \\ 0 \end{pmatrix}
                 + \beta\begin{pmatrix} 0 \\ 1 \end{pmatrix}}\\
  &\uncover<3->{= \begin{pmatrix} \mhl{\alpha} \\ \mhl*{\beta} \end{pmatrix}
     \begin{matrix} \rowlabel{salida $+z$} \\ \rowlabel{salida $-z$} \end{matrix}}
\end{align*}}%
\uncover<3->{$\alpha$ y $\beta$ son las \key{amplitudes}: dos números
\hl{complejos}, $\alpha,\beta\in\mathbb{C}$. No son probabilidades.}
\end{frame}

% ---------------------------------------------------------------------------
\begin{frame}{De amplitudes a probabilidades}{Normalización}
Si medimos con el imán en $z$:
\[
  P(0) = \explain{|\alpha|^2}{módulo al cuadrado}
  \qquad
  P(1) = |\beta|^2
\]
Para un número complejo, $|a + ib|^2 = a^2 + b^2$: siempre sale un número
real, nunca negativo.

El átomo siempre sale por algún lado, así que el estado está
\hl{normalizado}:
\keyeq{|\alpha|^2 + |\beta|^2 = 1}
\end{frame}

% ---------------------------------------------------------------------------
\begin{frame}{Un ejemplo}{Con números}
{\large
\[
  \ket{\psi} = \begin{pmatrix} 3/5 \\[0.5mm] 4i/5 \end{pmatrix}
\]}%
\begin{center}
\begin{tikzpicture}[x=1mm,y=1mm]
  \useasboundingbox (0,-41) rectangle (88,4);
  \node[anchor=base west, text=igMuted, font=\footnotesize] at (0,0) {salida};
  \foreach \x/\t in {39/amplitud, 58/{$|\cdot|^2$}, 78/probabilidad}
    \node[anchor=base, text=igMuted, font=\footnotesize] at (\x,0) {\t};
  \draw[igTrack, line width=0.5pt] (0,-2.5) -- (88,-2.5);
  % filas
  \node[anchor=base west, text=igFG] at (0,-12) {$+z \to \ket{0}$};
  \node[anchor=base, text=igFG] at (39,-12) {$\dfrac{3}{5}$};
  \node[anchor=base west, text=igFG] at (0,-24) {$-z \to \ket{1}$};
  \node[anchor=base, text=igFG] at (39,-24) {$\dfrac{4i}{5}$};
  % probabilidades (segundo paso)
  \node<2->[anchor=base, text=igFG] at (58,-12) {$\dfrac{9}{25}$};
  \node<2->[anchor=base, text=igFG] at (58,-24) {$\dfrac{16}{25}$};
  \node<2->[anchor=base, fill=igHL, draw=igTrack, rounded corners=1.2mm, text=igAccent,
            inner sep=1.2mm] at (78,-12) {36\,\%};
  \node<2->[anchor=base, fill=igHL, draw=igTrack, rounded corners=1.2mm, text=igAccent,
            inner sep=1.2mm] at (78,-24) {64\,\%};
  % barra: fracción de átomos en cada salida
  \fill<2->[igAccentB, rounded corners=0.8mm] (0,-35) rectangle (31.4,-31);
  \fill<2->[igAccent, rounded corners=0.8mm] (32.0,-35) rectangle (88,-31);
  \node<2->[anchor=north west, text=igAccentB, font=\footnotesize] at (0,-35.6) {$\ket{0}$: 36\,\%};
  \node<2->[anchor=north east, text=igAccent, font=\footnotesize] at (88,-35.6) {$\ket{1}$: 64\,\%};
\end{tikzpicture}
\end{center}
\uncover<2->{La $i$ no afecta a la probabilidad: $|4i/5|^2 = 16/25$.}
\end{frame}

% ---------------------------------------------------------------------------
\begin{frame}{Filtro $+z$, medimos $z$}{El post 02 en vectores}
\begin{center}
\begin{tikzpicture}[x=1mm,y=1mm,scale=1.3,transform shape]
  \draw[igAccentB, line width=0.8pt] (0,0) -- (8,0);
  \pic at (15,0) {sg={z/90/igAccentB}};
  \sgblock{15}
  \sgpass{15}{43}
  \pic at (43,0) {sg={z/90/igAccentB}};
  \sgout{43}{100\,\%}{0\,\%}
\end{tikzpicture}
\end{center}
Después del filtro, el estado es $\ket{0}$:
\[
  \ket{\psi} = \begin{pmatrix} 1 \\ 0 \end{pmatrix}
  \;\Longrightarrow\;
  P(0) = 1,\quad P(1) = 0 .
\]
La tabla de probabilidades dice \hl{100\,\% y 0\,\%}, lo mismo que vimos
con los imanes.
\end{frame}

% ---------------------------------------------------------------------------
\begin{frame}{Filtro $+x$, medimos $z$}{El post 02 en vectores}
\begin{center}
\begin{tikzpicture}[x=1mm,y=1mm,scale=1.15,transform shape]
  \draw[igAccentB, line width=0.8pt] (-2,0) -- (0,0);
  \pic at (7,0) {sg={z/90/igAccentB}};
  \sgblock{7}
  \sgpass{7}{31}
  \pic at (31,0) {sg={x/210/igAccent}};
  \sgblock{31}
  \sgpass{31}{55}
  \pic at (55,0) {sg={z/90/igAccentB}};
  \sgout{55}{50\,\%}{50\,\%}
\end{tikzpicture}
\end{center}
Lo que sale del filtro $+x$, escrito en la base $z$, es
\[
  \ket{\psi} = \begin{pmatrix} 1/\sqrt{2} \\[0.5mm] 1/\sqrt{2} \end{pmatrix}
  \;\Longrightarrow\;
  P(0) = P(1) = \Big|\frac{1}{\sqrt{2}}\Big|^2 = \frac12 .
\]
\hl{Mismo sistema, otro estado}: otra tabla de probabilidades.
\end{frame}

% ---------------------------------------------------------------------------
\begin{frame}{Del imán a las matemáticas}{Resumen}
\begin{center}
\begin{tikzpicture}[x=1mm,y=1mm]
  \node[anchor=base west, text=igMuted, font=\scriptsize] at (0,0) {\textls[120]{EXPERIMENTO}};
  \node[anchor=base west, text=igMuted, font=\scriptsize] at (55,0) {\textls[120]{MATEMÁTICAS}};
  \foreach \y/\l/\r in {
      -11/{Las dos salidas del imán $z$}/{La base $\ket{0}$, $\ket{1}$},
      -25/{Cómo preparamos el átomo}/{El estado $\begin{pmatrix}\alpha\\ \beta\end{pmatrix}$},
      -39/{Fracción de átomos en cada salida}/{$|\alpha|^2$ y $|\beta|^2$},
      -53/{Siempre sale por algún lado}/{$|\alpha|^2 + |\beta|^2 = 1$}}{
    \draw[igTrack, line width=0.5pt] (0,\y+6.5) -- (88,\y+6.5);
    \node[anchor=west, text width=44mm, text=igFG, align=left, inner sep=0,
          font={\hyphenpenalty10000\exhyphenpenalty10000}] at (0,\y) {\l};
    \node[text=igAccent] at (50,\y) {$\rightarrow$};
    \node[anchor=west, text=igFG, inner sep=0] at (55,\y) {\r};
  }
\end{tikzpicture}
\end{center}
\end{frame}

\closingslide{Un qubit, dos números}
  {Dos amplitudes complejas y normalizadas, en la base que eligió el imán.
   Guárdalo para el próximo post de la serie.}

\end{document}
